% ---------------------------------------------------------------- title
\begin{titlepage}
\begin{center}
\vspace*{0cm}
{\LARGE\bfseries EdgeGuard: Integrity Attestation and\\[0.2cm]
Layerwise Intrusion Detection for\\[0.2cm]
Quantized Neural Networks on Edge Accelerators\par}
\vspace{1.0cm}
{\large A Project Report\\[0.2cm] submitted in partial fulfilment of the requirements\\
for the award of the degree of\par}
\vspace{0.6cm}
{\Large\bfseries Master of Technology\par}
\vspace{0.3cm}
{\large in\\[0.2cm] Computer Science and Engineering\par}
\vspace{0.9cm}
{\large by\par}
\vspace{0.3cm}
{\Large\bfseries Arnav Deshpande\par}
{\large (CS24M513)\par}
\vspace{0.5cm}
{\normalsize Prepared with AI assistance from Claude (Anthropic);\\
see the Statement on the Use of AI Tools.\par}
\vfill
{\large\bfseries Department of Computer Science and Engineering\\[0.15cm]
Indian Institute of Technology Madras\\[0.15cm]
Chennai 600\,036\par}
\vspace{0.4cm}
{\large November 2026\par}
\end{center}
\end{titlepage}

% ---------------------------------------------------------------- certificate
\chapter*{Certificate}
\addcontentsline{toc}{chapter}{Certificate}
This is to certify that the report titled \textbf{EdgeGuard: Integrity Attestation and Layerwise
Intrusion Detection for Quantized Neural Networks on Edge Accelerators}, submitted by
\textbf{Arnav Deshpande (CS24M513)} to the Indian Institute of Technology Madras in partial
fulfilment of the requirements for the award of the degree of Master of Technology, is a
record of the project work carried out under my supervision.

\vspace{3cm}
\noindent
\begin{tabular}{@{}l@{\hspace{5cm}}l}
\rule{5cm}{0.4pt} & \\
\textbf{[Name of Guide]} & Place: Chennai \\
Project Guide & Date: \\
Department of Computer Science and Engineering & \\
IIT Madras &
\end{tabular}

% ---------------------------------------------------------------- AI statement
\chapter*{Statement on the Use of AI Tools}
\addcontentsline{toc}{chapter}{Statement on the Use of AI Tools}
This project was carried out with substantial assistance from Claude, a large language model
developed by Anthropic, which is credited as a co-author on the commits in the project
repository (\url{https://github.com/arnavd2001/edgeguard}). The assistance covered:

\begin{itemize}[leftmargin=*]
  \item drafting and debugging the Python implementation of the \textsc{EdgeGuard} framework,
        the attack simulators and the test suite;
  \item designing and running the experiments reported in \cref{ch:evaluation}, and generating
        the figures and tables from the recorded results;
  \item drafting the text of this report.
\end{itemize}

\noindent
Every experimental number in this report is produced by the scripts in
\code{experiments/} and is inserted into the document automatically from
\code{experiments/results/results.json}; no result was typed by hand. Every bibliographic
reference was checked against the publisher's page, the arXiv listing or a bibliographic
database (DBLP, Crossref, PMLR, ACL Anthology). An earlier draft of this report contained
results that had not been measured and references that could not be verified; those were
removed, and the corresponding sections were rewritten from measured data.

\noindent
The candidate is responsible for the content of this report, has reviewed it, and should be
able to explain every design decision and result in it.

\vspace{1.5cm}
\noindent Arnav Deshpande \hfill Date: \underline{\hspace{3cm}}

% ---------------------------------------------------------------- acknowledgements
\chapter*{Acknowledgements}
\addcontentsline{toc}{chapter}{Acknowledgements}
\textit{[To be written by the candidate.]}

% ---------------------------------------------------------------- abstract
\chapter*{Abstract}
\addcontentsline{toc}{chapter}{Abstract}

Neural networks deployed on edge accelerators are stored and executed as 8-bit integer (INT8)
tensors in memory that an attacker may be able to read or corrupt. Three families of attack
are well documented: bit flips induced in DRAM that destroy or hijack a model, adversarial
inputs that steer it, and extraction of weights or memorised training data from it. Defences
in the literature usually address one family at a time. This report describes
\textsc{EdgeGuard}, a Python framework that combines two independent checks and a rule-based
classifier that labels what it sees.

The first check is a Merkle tree of SHA-256 digests over the model's weight tensors, with an
ECDSA P-256 signature on the root. It detects any change to the stored weights and names the
tensor that changed. The second check is a per-layer activation monitor calibrated with
Welford's streaming algorithm directly on INT8 activations, scoring each inference with a
channel-wise $z$-score and a Mahalanobis distance. A classifier maps the joint outcome to one of
eight reachable threat labels and emits a layer-by-layer explanation. The framework includes
simulators for bit-flip, supply-chain, adversarial-input, training-data-extraction and combined
(``synergy'') attacks, and a generator of INT8 weight sets shaped like five open-weight
language models.

Evaluation is carried out on simulated activations, with \NPerScenario{} inputs per scenario
across \NSeeds{} seeds. The Merkle check flagged every weight-modifying input in every scenario
and localised the modified tensor exactly in every case. The activation monitor alone flagged
only \MonSynergy{} of synergy inputs, which perturb weights and suppress the activation signal,
so the combined pipeline catches inputs that activation monitoring alone misses. A fixed
$4\sigma$ threshold produced a \EmpFPRZFour{} false-positive rate on clean inputs, which matches
the \PredFPRZFour{} predicted by a multiple-comparisons argument over \MChannels{} monitored
channels; a Bonferroni-corrected threshold of $z^*=\ZBonf$ reduced it to \EmpFPRBonf{} while
still flagging all three activation-level attack scenarios. Re-hashing a 1.4\,GiB weight set took
\MerkVerLlamaPB{}\,ms on the test machine, which rules out full attestation on every inference
for language-model-sized weights and motivates the amortised schemes discussed in
\cref{ch:discussion}. Because the attack generators and the detector were written together,
the activation-level results measure the internal consistency of the simulation rather than
robustness to an adaptive attacker; this limitation and a plan to evaluate on real models on a
Qualcomm Hexagon NPU are set out in \cref{ch:conclusion}.

\vspace{0.5cm}
\noindent\textbf{Keywords:} edge AI security, quantized neural networks, bit-flip attacks,
Merkle trees, anomaly detection, Mahalanobis distance, training-data extraction, NPU.

% ---------------------------------------------------------------- lists
\cleardoublepage
\tableofcontents
\cleardoublepage
\listoffigures
\addcontentsline{toc}{chapter}{List of Figures}
\cleardoublepage
\listoftables
\addcontentsline{toc}{chapter}{List of Tables}

% ---------------------------------------------------------------- abbreviations
\chapter*{Abbreviations}
\addcontentsline{toc}{chapter}{Abbreviations}
\begin{tabular}{@{}p{2.6cm}p{11cm}@{}}
AUC   & Area under the (ROC) curve \\
BFA   & Bit-flip attack \\
CVD   & Colour-vision deficiency \\
DMA   & Direct memory access \\
DRAM  & Dynamic random-access memory \\
ECDSA & Elliptic Curve Digital Signature Algorithm \\
EM    & Electromagnetic (side channel) \\
FFN   & Feed-forward network (transformer MLP block) \\
FPR   & False-positive rate \\
FP32  & 32-bit IEEE-754 floating point \\
HTP   & Hexagon Tensor Processor \\
INT8  & 8-bit signed two's-complement integer \\
LLM   & Large language model \\
LSB / MSB & Least / most significant bit \\
NPU   & Neural processing unit \\
PBFA  & Progressive bit-flip attack \\
PTQ   & Post-training quantization \\
QAT   & Quantization-aware training \\
QNN   & Qualcomm AI Engine Direct SDK (``Qualcomm Neural Network'') \\
ROC   & Receiver operating characteristic \\
SHA-256 & Secure Hash Algorithm, 256-bit output \\
TEE   & Trusted execution environment \\
TOCTOU & Time of check to time of use \\
TPR   & True-positive rate \\
\end{tabular}

% ---------------------------------------------------------------- notation
\chapter*{Notation}
\addcontentsline{toc}{chapter}{Notation}
\begin{tabular}{@{}p{2.6cm}p{11cm}@{}}
$L_i$                     & $i$-th weight tensor of a model, $i = 0,\dots,n-1$ \\
$n$                       & number of weight tensors (Merkle leaves) \\
$\Hash(\cdot)$            & SHA-256 \\
$h_i$                     & leaf digest of $L_i$ \\
$r$                       & Merkle root \\
$\sigma$                  & ECDSA signature on $r$ \\
$\mathbf{a}_\ell\in\Z^{d}$ & INT8 activation vector at monitored layer $\ell$, dimension $d$ \\
$\boldsymbol\mu_\ell,\boldsymbol\sigma_\ell$ & per-channel calibration mean and standard deviation \\
$\hat\Sigma_\ell$         & calibration covariance; $\hat\Sigma_\ell^{+}$ its pseudo-inverse \\
$Z_\ell$                  & max channel $|z|$-score at layer $\ell$ \\
$D_\ell$                  & Mahalanobis distance at layer $\ell$ \\
$\tau_z,\tau_D$           & $z$-score and Mahalanobis thresholds \\
$m$                       & total monitored channels, $m=\sum_\ell d_\ell$ \\
$s, z_0$                  & quantization scale and zero point \\
\end{tabular}
